\chapter{Aspects of Measurement Data Analysis}

In current scientific research, the analysis of measurement data is crucial for drawing meaningful conclusions. More of one measurement, always under the same conditions, is called a \textit{series of measurements}. The measurement result is the mean value of the series of measurements. There are two types of measurements:
\begin{itemize}
    \item Direct Measurements: The quantity being measured is directly obtained from the measurement process.
    \item Indirect Measurements: The quantity being measured is calculated from other directly measured quantities using a mathematical relationship.
\end{itemize}

However, the accuracy of measurements can be affected by various types of errors. These are:
\begin{description}
    \item[Gross errors] These are errors in the strict sense, often resulting from incorrect behavior on the part of the observer or damaged, unusable measuring instruments. Gross errors are always avoidable.
    \item[Systematic deviations] Also known as systematic errors, these arise from inaccurate measurement methods or faulty measuring instruments, such as calibration or parallax errors. Systematic deviations can be reduced to a negligible extent through care and high-quality instruments. Correction is possible if the deviations can be readily identified theoretically or empirically.
    \item[Random deviations] Also known as random or statistical errors, these are fluctuations within a series of measurements when the measurement is repeated. These errors are then analyzed using probability theory and statistics.
\end{description}

In the analysis of measurement data, it is essential to consider how the error has been spread. A regression analysis can be used to determine the relationship between the measured quantity and other variables.

\section{Mathmatical Statistics}

When given a measure, the measurement error should also be provided. There are two options for expressing the measurement error:
\begin{itemize}
    \item Absolute error: This is the difference between the measured value and the expected true value. It is expressed in the same units as the measured quantity.
    \item Relative error: This is the ratio of the absolute error to the measured value, often expressed as a percentage.
\end{itemize}

Each of the measurements in a series of measurements is a random variable which follows a certain probability distribution. Some important values of that random variable $X$ are:
\begin{itemize}
    \item The mean value $\mu$ is the expected value of the random variable, calculated as $\mu = E[X] =\langle X \rangle$.
    \begin{equation*}
        \langle X \rangle = \dfrac{1}{N}\sum_{i=1}^N X_i
    \end{equation*}
    \item The variance $\sigma^2$ is a measure of the spread of the random variable, calculated as $\sigma^2 = Var[X]$.
    \begin{equation*}
        Var[X] = \dfrac{1}{N-1}\sum_{i=1}^N (X_i - \langle X \rangle)^2
    \end{equation*}
    It is important to note that the variance is calculated using $N-1$ in the denominator, which is known as Bessel's correction. This correction is applied to provide an unbiased estimate of the population variance when using a sample of data.
    \item The standard deviation $\sigma$ represents the average amount of variability in the data.
    \begin{equation*}
        \sigma = \sqrt{Var[X]}
    \end{equation*}
    When the number of measurements $N$ increases, the variance and standard deviation tend to be constant.

    \item The standard error of the mean (SEM) is a measure of how much the sample mean is expected to vary from the true population mean. It is calculated as:
    \begin{equation*}
        SEM = \dfrac{\sigma}{\sqrt{N}}
    \end{equation*}
    When the number of measurements $N$ increases, the standard error of the mean decreases, indicating that the sample mean is likely to be closer to the true population mean. In addition, given to the square root in the denominator, the standard error of the mean decreases at a slower rate than the standard deviation as $N$ increases.
    \begin{itemize}
        \item In order to improve the accuracy of the measurement by $3$ times, the number of measurements must be increased by $3^2 = 9$ times.
        \item In order to improve the accuracy of the measurement by $10$ times, the number of measurements must be increased by $10^2 = 100$ times.
    \end{itemize}
    As they reader can see, there is a point where increasing the number of measurements does not significantly improve the accuracy of the measurement.
\end{itemize}

It has empirically been found that the distribution of measurement errors, specially when the number of measurements is large, can be approximated by a normal distribution, also known as a Gaussian distribution. It has the following interesing properties:
\begin{itemize}
    \item The mean, median, and mode of a normal distribution are all equal.
    \item The normal distribution is symmetric around its mean, meaning that the left and right sides of the distribution are mirror images of each other.
    \item The normal distribution is defined by two parameters: the mean $\mu$ and the standard deviation $\sigma$. The mean determines the center of the distribution, while the standard deviation determines the spread of the distribution.
    \item The closer a value is to the mean, the higher its probability density. In fact:
    \begin{align*}
        P[|X - \mu| < \sigma] &\approx 68.27\% \\
        P[|X - \mu| < 2\sigma] &\approx 95.45\% \\
        P[|X - \mu| < 3\sigma] &\approx 99.73\%
    \end{align*}
\end{itemize}


\section{Error Propagation}

In indirect measurements, the quantity being measured is calculated from other directly measured quantities using a mathematical relationship. Before moving on to quantities that depend on several measured quantities, we will first consider the case of a quantity $f$ that depends on a single measured quantity $x$. The error in $f$ can be approximated using the total differential of $f$ with respect to $x$:
\begin{equation*}
    \Delta f = \left| \dfrac{df}{dx} \right| \Delta x
\end{equation*}

However, there are cases where the quantity $f$ depends on several measured quantities $x_1, x_2, \ldots, x_n$. In such cases, the error in $f$ just because of the errors in one of the measured quantities $x_i$ can be approximated using the partial derivative of $f$ with respect to $x_i$:
\begin{equation*}
    (\Delta f)_i = \left| \dfrac{\partial f}{\partial x_i} \right| \Delta x_i
\end{equation*}

However, when considering the total error in $f$ due to the errors in all the measured quantities, there are two common approaches:
\begin{itemize}
    \item The maximum error approach, which simply sums up the individual errors from each measured quantity:
    \begin{equation*}
        (\Delta f)_{\max} = \sum_{i=1}^n (\Delta f)_i = \sum_{i=1}^n \left| \dfrac{\partial f}{\partial x_i} \right| \Delta x_i
    \end{equation*}
    \item Error measurements are often assumed to be independent, and from an statistical point of view, they partially cancel each other out. In this case, the Gaussian error propagation formula is used, which combines the individual errors in quadrature:
    \begin{equation*}
        (\Delta f)_{\text{Gaussian}} = \sqrt{\sum_{i=1}^n \left(\Delta f\right)_i^2} = \sqrt{\sum_{i=1}^n \left( \left| \dfrac{\partial f}{\partial x_i} \right|^2\cdot \Delta x_i^2 \right)}
    \end{equation*}
\end{itemize}

\section{Regression Analysis}

Regression analysis is a statistical method used to model the relationship between a dependent variable and one or more independent variables. Let us name the dependent variable as $y$ and the independent variables as $x\in \mathbb{R}^{\ell}$, where $\ell$ is the number of independent variables. Lets say we have a set of $N$ measurements of the dependent variable, denoted as $y_1, y_2, \ldots, y_N$, corresponding to the independent variable values $x_1, x_2, \ldots, x_N$ (taking into account that the independent variable can be a vector). The goal of regression analysis is to find a function $f(x)$ that best fits the data points $(x_i, y_i)$ for $i = 1, 2, \ldots, N$. That function has $n$ parameters, which we will denote as $a_1, a_2, \ldots, a_n$. Then, our model is given by:
\begin{equation*}
    f: \mathbb{R}^{\ell} \times \mathbb{R}^n
\end{equation*}

It should be noted that our aim is to find the parameters $a_1, a_2, \ldots, a_n$ so that the function $f$ depends only on the independent variable $x$.

\subsection{Linear Regression}

Linear regression is a specific type of regression analysis where the model function $f$ is linear in the parameters. This means that $f$ can be expressed as:
\begin{equation*}
    f(x, a) = \Phi(x) \cdot a = \sum_{j=1}^n \Phi_j(x) a_j
\end{equation*}
where $\Phi: \mathbb{R}^{\ell} \to \mathbb{R}^n$ is a vector of basis functions such thtat $\Phi_1(x)=1$, and $a$ is the vector of parameters. It should be noted that the basis functions $\Phi_j(x)$ can be nonlinear functions of the independent variable $x$, but the model is still considered linear because it is linear in the parameters $a_j$.\\

Ideally, we would like to find the parameters $a$ such that the model function $f(x_i, a)$ perfectly fits the data points $y_i$ for all $i$. For each data point, we can stablish the following equation:
\begin{equation*}
    y_i = f(x_i, a) = \Phi(x_i) \cdot a
\end{equation*}
We will then have a system of $N$ equations (one for each data point) with $n$ unknowns (the parameters $a_j$). That system can be written in matrix form as:
\begin{equation*}
    K a = y
\end{equation*}
where $K$ is the $N \times n$ matrix whose $i$-th row is given by $\Phi(x_i)$, and $y$ is the vector of measurements. However, given that having more data points improves the accuracy of the measurement, it is common to have $N \gg n$. In this case, it is not possible to find a set of parameters $a$ that perfectly fits all the data points:
\begin{equation*}
    K a - y \neq 0\qquad \forall a\in \mathbb{R}^n
\end{equation*}

Our goal is then to find the parameters $a$ that minimize the error between the model predictions and the actual measurements, i.e., we want to find $a^\ast\in \mathbb{R}^n$ such that:
\begin{equation*}
    \| K a^\ast - y \| = \min_{a\in \mathbb{R}^n} \| K a - y \|
\end{equation*}

In order to get rid of the square root, our problem can be equivalently written as finding $a^\ast\in \mathbb{R}^n$ such that:
\begin{equation*}
    \| K a^\ast - y \|^2 = \min_{a\in \mathbb{R}^n} \| K a - y \|^2
\end{equation*}

We have just described the \emph{least squares} problem, which is a common method for performing linear regression. Analiticaly, the solution to the least squares problem can be found by finding the minimum of the function:
\Func{F}{\mathbb{R}^n}{\mathbb{R}}{a}{\| K a - y \|^2=\sum_{i=1}^N (f(x_i, a) - y_i)^2}

That can be done by ensuring that $\nabla F(a) = 0$, which has $n$ equations (each of the partial derivatives must be zero) with $n$ unknowns (the parameters $a_j$).
\begin{observacion}
    In MATLAB, the least squares problem can be solved using the backslash operator, which is a convenient way to solve linear systems of equations. Given a matrix $K$ and a vector $y$, you can find the parameters $a$ that minimize the least squares error by using \texttt{a = K\textbackslash y}
\end{observacion}

\subsubsection{Deciding which model fits better}

When we have several models to choose from, it may be difficult to decide which one fits the data better. One common approach is to use the coefficient of determination, denoted as $R^2$, which measures the proportion of the variance in the dependent variable that is predictable from the independent variables.
\begin{equation*}
    R^2 = \dfrac{\sum_{i=1}^N (f(x_i, a) - \langle y \rangle)^2}{\sum_{i=1}^N (y_i - \langle y \rangle)^2}
\end{equation*}
where $\langle y \rangle$ is the mean of the observed data $y_i$. The value of $R^2$ ranges from 0 to 1, where a value closer to 1 indicates a better fit of the model to the data.
\begin{observacion}
    It should be noted that a high $R^2$ only informs that the model fits the data \emph{we currently have} well, but it does not necessarily mean that the model will predict future data accurately\footnote{``With enough parameters, I can fit an elephant'' - John von Neumann}. Therefore, it is important to consider other factors, such as the complexity of the model and the underlying physical principles, when deciding which model is more appropriate for a given set of data.
\end{observacion}

\subsection{Nonlinear Regression}

Nonlinear regression is a type of regression analysis where the model function $f$ is nonlinear in the parameters. This means that $f$ cannot be expressed as a linear combination of the parameters. Nonlinear regression models can take various forms, but the more relevants are the exponential and logarithmic models. These models can be linearized by applying specific transformations to the data, allowing us to use linear regression techniques to estimate the parameters. For example, consider the exponential model:
\begin{equation*}
    f(x, a) = a_1 e^{a_2 x}a
\end{equation*}
By taking the natural logarithm of both sides, we can linearize the model:
\begin{equation*}
    \ln(f(x, a)) = \ln(a_1) + a_2 x
\end{equation*}
This transformation allows us to apply linear regression techniques to estimate the parameters $\ln(a_1)$ and $a_2$. We can then exponentiate the estimated value of $\ln(a_1)$ to obtain the estimate for $a_1$.